\begin{proposition}[3.2.6]
\label{IV.2.3.2.6-fr}
Sean $X$ un preesquema localmente noetheriano y $\mathcal F$ un
$\mathcal O_X$-módulo coherente. Entonces existe una descomposición
irredundante reducida
$(\mathcal F^{(x)})_{x\in\operatorname{Ass}(\mathcal F)}$, constituida por
$\mathcal O_X$-módulos coherentes tales que, para todo
$x\in\operatorname{Ass}(\mathcal F)$, se tiene
$\operatorname{Ass}(\mathcal F^{(x)})=\{x\}$. Para todo
$x\in\operatorname{Ass}(\mathcal F)$ tal que $\overline{\{x\}}$ no sea
inmerso, $\mathcal F^{(x)}$ queda determinado de modo único como la imagen del
homomorfismo canónico $\mathcal F\to i_*(i^*(\mathcal F))$, donde $i$ es el
morfismo canónico $\operatorname{Spec}(\boldsymbol{k}(x))\to X$.
\end{proposition}

Para todo $x\in\operatorname{Ass}(\mathcal F)$, sea $U$ una vecindad abierta
afín de $x$, con anillo $A$, y sea $\mathcal F|U=\widetilde M$, donde $M$ es
un $A$-módulo de tipo finito. Se sabe (Bourbaki, \emph{Alg. comm.}, cap.~IV,
\S~1, núm.~1, prop.~4) que existe un submódulo $N$ de $M$ tal que, si se pone
$P=M/N$, se tiene $\operatorname{Ass}(P)=\{x\}$ y
$\operatorname{Ass}(N)=\operatorname{Ass}(M)\setmin\{x\}$. Sea
$\mathcal G=\widetilde P$, que es un $\mathcal O_U$-módulo cuasicoherente, y
sea $j$ la inyección canónica $U\to X$; sea
$u:j^*(\mathcal F)\to\mathcal G$ el homomorfismo sobreyectivo correspondiente
al homomorfismo $M\to P$; de él se deduce un homomorfismo
$j_*(u):j_*(j^*(\mathcal F))\to j_*(\mathcal G)$ y, por composición, un
homomorfismo
\[
v:\mathcal F\xrightarrow{\rho_{\mathcal F}}j_*(j^*(\mathcal F))
\xrightarrow{j_*(u)}j_*(\mathcal G)
\]
del cual $u$ es la restricción a $U$; denotaremos por
$\mathcal F^{(x)}$ la imagen de $\mathcal F$ por este homomorfismo, que es un
$\mathcal O_X$-módulo coherente
(\textbf I, 6.1.1). Se tiene $\operatorname{Ass}(j_*(\mathcal G))=\{x\}$ en virtud de
(3.1.13) y, \emph{a fortiori} (3.1.7),
$\operatorname{Ass}(\mathcal F^{(x)})=\{x\}$, puesto que
$\mathcal F^{(x)}\ne0$. Además, si

\oldpage[IV-2]{43}
$\mathcal N^{(x)}=\operatorname{Ker}(v)$, se tiene
$\mathcal N^{(x)}|U=\widetilde N$, y por tanto
$x\notin\operatorname{Ass}(\mathcal N^{(x)})$. De ello resulta que el
homomorfismo
$\mathcal F\to\bigoplus_{x\in\operatorname{Ass}(\mathcal F)}
\mathcal F^{(x)}$ es inyectivo, pues su núcleo $\mathcal H$ está contenido en
todos los $\mathcal N^{(x)}$, de modo que $\operatorname{Ass}(\mathcal H)$
está contenido en la intersección de los
$\operatorname{Ass}(\mathcal N^{(x)})$, que es vacía; por consiguiente (3.1.5),
$\mathcal H=0$. Teniendo en cuenta que
$\operatorname{Ass}(\mathcal F)$ es localmente finito (3.1.6), es claro que
$(\mathcal F^{(x)})_{x\in\operatorname{Ass}(\mathcal F)}$ es una
descomposición irredundante reducida de $\mathcal F$ que satisface las
condiciones del enunciado. La caracterización de $\mathcal F^{(x)}$ cuando
$\overline{\{x\}}$ no es inmerso resulta de Bourbaki,
\emph{Alg. comm.}, cap.~IV, \S~2, núm.~3, prop.~5, puesto que la cuestión es
local, y teniendo en cuenta (\textbf I, 1.6.7).
